\documentclass[11pt]{article}
\usepackage[margin=1in]{geometry}
\usepackage{amsmath}
\usepackage{booktabs}
\usepackage{longtable}
\usepackage{enumitem}
\usepackage{xcolor}
\usepackage{hyperref}
\usepackage{listings}

\lstset{
    basicstyle=\small\ttfamily,
    breaklines=true,
    frame=single
}

\title{Codebase Organization Guide\\
\large Singum Traffic Flow Research}
\author{}
\date{\today}

\begin{document}
\maketitle

\tableofcontents
\newpage

%==============================================================================
\section{Overview}
%==============================================================================

This codebase implements the singum model for traffic flow prediction as described in the PNAS paper ``Homogenization of a two-dimensional network for flow prediction and network routing.'' The code:

\begin{itemize}[noitemsep]
    \item Homogenizes a road network into a continuum using Voronoi cells (singums)
    \item Computes effective conductivity tensors $k_{ij}$ (Paper Eqs. 20--21)
    \item Solves the anisotropic diffusion PDE $\nabla \cdot (K \nabla V) = 0$ using FEM (Paper Eq. 11)
    \item Performs greedy routing on the potential field
    \item Generates Figures 5c and 5d for the paper
\end{itemize}

%==============================================================================
\section{Directory Structure}
%==============================================================================

\begin{verbatim}
singum_research/
|-- case2.py                  # Main pipeline (Fig 5c, 5d) with FEM solver
|-- singum.py                 # Voronoi cell generation
|-- load_data.py              # Data loading utilities
|-- symmetric_transportation.py
|-- README.md                 # Quick start instructions
|-- requirements.txt          # Python dependencies
|-- .gitignore
|
|-- data/                     # Input data (git-ignored)
|   |-- filtered_roads.geojson
|   |-- manhattan_boundary.geojson
|   |-- nyc_boundary.geojson
|   |-- singum_voronoi.geojson
|   +-- lion/                 # LION geodatabase
|
|-- output/                   # Generated files (git-ignored)
|
|-- docs/                     # Documentation
|   |-- case2_functions_documentation.tex/pdf
|   +-- codebase_organization_guide.tex/pdf
|
+-- archive/                  # Old reference files (not used)
\end{verbatim}

%==============================================================================
\section{File Descriptions}
%==============================================================================

\subsection{Active Source Files}

\begin{longtable}{p{4.5cm}p{9cm}}
\toprule
\textbf{File} & \textbf{Description} \\
\midrule
\endfirsthead

\texttt{case2.py} & Main pipeline script. Implements:
\begin{itemize}[noitemsep,topsep=0pt]
    \item FEM solver for anisotropic diffusion PDE (Eq. 11)
    \item K tensor computation with $k^2$ correction (Eqs. 20--21)
    \item Greedy and Dijkstra routing
    \item Figure generation (5c, 5d)
\end{itemize} \\
\midrule
\texttt{singum.py} & Generates Voronoi cells (singums) from road network nodes. Outputs \texttt{singum\_voronoi.geojson}. Run once to create tessellation. \\
\midrule
\texttt{load\_data.py} & Utility functions to load LION geodatabase and NYC boundary data. Used by \texttt{singum.py}. \\
\midrule
\texttt{symmetric\_transportation.py} & Symmetric transportation model implementation. \\
\bottomrule
\end{longtable}

\subsection{Data Files (in \texttt{data/})}

\begin{longtable}{p{4.5cm}p{9cm}}
\toprule
\textbf{File} & \textbf{Description} \\
\midrule
\endfirsthead

\texttt{filtered\_roads.geojson} & NYC road network (55 MB). Contains road segments with NodeIDFrom, NodeIDTo, TrafDir. \\
\midrule
\texttt{manhattan\_boundary.geojson} & Manhattan island polygon boundary. \\
\midrule
\texttt{nyc\_boundary.geojson} & All NYC borough boundaries. Used to extract Manhattan. \\
\midrule
\texttt{singum\_voronoi.geojson} & Voronoi cells generated by \texttt{singum.py}. Each cell corresponds to a road network node (singum particle). \\
\midrule
\texttt{lion/} & NYC LION geodatabase. Source data for road network. \\
\bottomrule
\end{longtable}

%==============================================================================
\section{Output Files (Generated by Pipeline)}
%==============================================================================

When you run \texttt{case2.py}, these files are created in \texttt{output/}:

\begin{longtable}{p{5.5cm}p{8cm}}
\toprule
\textbf{File} & \textbf{Generated By} \\
\midrule
\endfirsthead

\multicolumn{2}{l}{\textbf{Intermediate Files}} \\
\midrule
\texttt{filtered\_roads\_2263.geojson} & \texttt{align\_and\_clip()} \\
\texttt{manhattan\_boundary\_2263.geojson} & \texttt{align\_and\_clip()} \\
\texttt{filtered\_roads\_manhattan\_2263.geojson} & \texttt{align\_and\_clip()} \\
\texttt{graph\_edges.txt} & \texttt{build\_graph\_from\_roads()} \\
\texttt{nodes\_points\_2263.geojson} & \texttt{build\_graph\_from\_roads()} \\
\texttt{nodes\_with\_phi\_2263.geojson} & \texttt{solve\_node\_potential()} \\
\midrule

\multicolumn{2}{l}{\textbf{Voronoi/K Tensor Outputs}} \\
\midrule
\texttt{singum\_voronoi\_with\_K.geojson} & \texttt{build\_K\_from\_voronoi()} -- includes $k^1 + k^2$ \\
\texttt{singum\_voronoi.png} & \texttt{build\_K\_from\_voronoi()} \\
\texttt{K\_centroids\_2263.csv} & \texttt{build\_K\_from\_voronoi()} \\
\midrule

\multicolumn{2}{l}{\textbf{Continuum/Figure Outputs (FEM)}} \\
\midrule
\texttt{manhattan\_continuum\_V.png} & \texttt{solve\_continuum\_and\_plot()} \\
\texttt{fig5c\_nba\_*.png} & \texttt{case\_5c\_nba\_game()} \\
\texttt{fig5c\_nba\_arrival\_*.png} & \texttt{case\_5c\_nba\_arrival()} -- uses FEM \\
\texttt{fig5d\_gw\_commute.png} & \texttt{case\_5d\_gw\_commute()} \\
\bottomrule
\end{longtable}

%==============================================================================
\section{How to Run}
%==============================================================================

\subsection{Prerequisites}

Install Python dependencies:
\begin{lstlisting}[language=bash]
pip install numpy pandas geopandas matplotlib scipy shapely fiona rasterstats
\end{lstlisting}

Or use the requirements file:
\begin{lstlisting}[language=bash]
pip install -r requirements.txt
\end{lstlisting}

\subsection{Generate Voronoi Cells (One Time)}

If \texttt{data/singum\_voronoi.geojson} doesn't exist:
\begin{lstlisting}[language=bash]
python singum.py
\end{lstlisting}

\subsection{Run Main Pipeline}

\begin{lstlisting}[language=bash]
python case2.py
\end{lstlisting}

This runs the full pipeline and generates all figures using:
\begin{itemize}[noitemsep]
    \item FEM solver with P1 triangular elements ($\sim$163k nodes, $\sim$327k triangles)
    \item K tensor with $k^2$ correction for non-symmetric networks
\end{itemize}

\subsection{Run Individual Figures}

In Python:
\begin{lstlisting}[language=Python]
from case2 import run_fig5c, run_fig5c_arrival, run_fig5d

run_fig5c()           # NBA dispersal from MSG
run_fig5c_arrival()   # Arrival at MSG (uses FEM)
run_fig5d()           # GW Bridge commute
\end{lstlisting}

%==============================================================================
\section{Implementation Notes}
%==============================================================================

\subsection{FEM Solver}

The continuum PDE $\nabla \cdot (K \nabla V) = 0$ is solved using:
\begin{itemize}[noitemsep]
    \item P1 (linear) triangular finite elements
    \item Delaunay triangulation of interior grid points
    \item Dirichlet boundary conditions via row replacement
    \item \texttt{scipy.sparse.linalg.spsolve} for linear system
\end{itemize}

\subsection{K Tensor Computation}

The effective conductivity follows Paper Eqs. 20--21:
\begin{equation}
    k_{ij} = k^1_{ij} + k^2_{ij}
\end{equation}
where $k^1$ is the Cauchy-Born rule term and $k^2$ is the correction for non-central-symmetric networks (like Manhattan's irregular road grid).

%==============================================================================
\section{Git Ignore Rules}
%==============================================================================

The \texttt{.gitignore} excludes:

\begin{itemize}[noitemsep]
    \item \texttt{data/} -- Large input files (55+ MB)
    \item \texttt{output/} -- All generated files
    \item \texttt{*.geojson} -- GeoJSON files anywhere
    \item \texttt{*.png} -- Generated images
    \item \texttt{*.csv} -- Generated data
    \item \texttt{\_\_pycache\_\_/} -- Python cache
\end{itemize}

Documentation PDFs in \texttt{docs/} are \textbf{not} ignored.

%==============================================================================
\section{Dependencies Between Files}
%==============================================================================

\begin{verbatim}
singum.py
  |-- imports from: load_data.py
  |-- reads: data/lion/lion.gdb (via load_data)
  |-- reads: data/manhattan_boundary.geojson
  +-- writes: data/singum_voronoi.geojson

case2.py
  |-- reads: data/filtered_roads.geojson
  |-- reads: data/manhattan_boundary.geojson
  |-- reads: data/singum_voronoi.geojson
  |-- FEM solver: solve_fem_anisotropic()
  |     |-> _create_fem_mesh() (Delaunay triangulation)
  |     |-> _assemble_fem_system() (stiffness matrix)
  |     +-> spsolve() (linear solve)
  |-- K tensor: build_K_from_voronoi()
  |     +-> computes k^1 + k^2 (Eqs. 20-21)
  +-- writes: output/*

load_data.py
  |-- reads: data/lion/lion.gdb
  +-- reads: data/nyc_boundary.geojson
\end{verbatim}

\end{document}
